% v0.11-I Introduction text retained; only author TODO markers/comments added.
% Replace sections/introduction.tex when applying this review patch.
\providecommand{\dtctodo}[2]{\textbf{[TODO #1: #2]}}
\section{Introduction}
GPU L1 caches need storage and service bandwidth for new requests and outstanding accesses~\cite{kroft,hierarchical_mha,bandwidth}.
In the reservation-based L1 studied here, a miss reserves a way until its response returns~\cite{li2015}.
When every way in a set is reserved, another same-set miss stalls even if miss-tracking entries remain free.
A larger cache can ease this constraint, but retains the same tag-to-data allocation rule.

Decoupled-Tag Cache (DTC) separates address lookup from data retention~\cite{zhao2024}.
New requests search tag mappings.
Admitted transactions and their consumers keep physical-line indices.
DTC can therefore replace a mapping while preserving the old transaction's destination and its consumers' data.
Returned data stay in a per-SM pool.
The pending-instruction buffer (PIB) stores indices and dependencies instead of complete result vectors.
Consumers later collect their data from the pool.

In-order DTC (IO) completes instructions in admission order.
It releases replaced lines at the replacing instruction's completion.
A ready instruction retains its PIB entry while older instructions wait.
Out-of-order DTC (OO) tracks consumer references and permits eligible instructions to complete independently.
It can release replaced lines sooner, but requires reference state, waiter updates, and ready selection.

Tag replacement introduces a separate cost.
Once a mapping is replaced, later requests cannot hit or merge with its old allocation.
The data may still occupy the pool to serve existing consumers.
Collection and reference updates also require service bandwidth~\cite{dtc-artifact}.
We ask when this organization is preferable to a larger searchable cache or conventional waiting buffer.
% These three IDs mark the specific argument that remains open in paragraph 4.
% Correct separation of lifetimes does not yet establish superior performance or lower cost.
\dtctodo{COV-01}{Tag-coverage evidence}
\dtctodo{CMP-01}{Performance attribution versus C+}
\dtctodo{HW-01}{Realize and cost finite services}

The IO/OO designs originate in the Unified-Cache GPGPU work~\cite{zhao2024}.
This study adds a 64-SM model with finite lookup, refill, wakeup, collection, and retirement services~\cite{dtc-artifact}.
Comparisons include cache expansion and larger conventional waiting buffers.
A separate control matches IO/OO waiting depth.
With 128 PIB entries, OO is $1.963\times$ faster than IO on BICG.
It still reaches only $0.533\times$ the performance of the enhanced conventional design, \cplus{}.
On GEMM, OO instead reaches $1.097\times$ the performance of \cplus{}.
These results distinguish faster consumer completion from an advantage over conventional cache expansion.
% TODO NOV-01: Finalize the independent contribution relative to the thesis and nearest designs.
% Keep original design attribution. An account of validation work is not itself a new mechanism.
\dtctodo{NOV-01}{Independent contribution and prior-art boundary}
